%%%%%%%%%%%%%%%%%%%%%%%%%%%%%%%%%%%%%%%%%%%%%%%%%%%%%%%%%%%%%%%%%%%%%%
% Lớp 10 | Hệ thức lượng trong tam giác | Giá trị lượng giác của góc (0 −180∘)
% Dạng 0H4?1-3: Biến đổi, rút gọn biểu thức lượng giác
% Số câu: 23  (sinh tự động bởi tach_cau.py; từ nay thầy sửa trực tiếp file này)
%%%%%%%%%%%%%%%%%%%%%%%%%%%%%%%%%%%%%%%%%%%%%%%%%%%%%%%%%%%%%%%%%%%%%%

%%%%%%%%%%%%%%%% TRẮC NGHIỆM 4 PHƯƠNG ÁN %%%%%%%%%%%%%%%%

% nguồn: MaHoaNH_TracNghiem_Lop10
\begin{ex}%[0H4H1-3]%[ma:f209b46e]
	Cho góc  $\alpha$ thỏa mãn $90^{\circ}<\alpha<180^{\circ}$ và $\sin\alpha=\dfrac{1}{3}$. Tính giá trị của biểu thức \break$A=\sin\alpha\cdot\cot\alpha+\cos\alpha\cdot\tan\alpha$.
	\choice {\True $\dfrac{1-2\sqrt{2}}{3}$}{ $\dfrac{1+2\sqrt{2}}{3}$}{  $\dfrac{2\sqrt{2}}{3}$}{ $\dfrac{-\sqrt{2}+3}{3}$}
	 \loigiai{
		Ta có $A=\sin\alpha\cdot\dfrac{\cos\alpha}{\sin\alpha}+\cos\alpha\cdot\dfrac{\sin\alpha}{\cos\alpha}=\cos\alpha+\sin\alpha$.\\ 
		Từ $\sin^{2}\alpha+\cos^{2}\alpha=1  \Rightarrow \cos^{2}\alpha=1-\sin^{2}\alpha=1-\left(\dfrac{1}{3}\right)^{2}=\dfrac{8}{9}$.\\
		Vì $90^{\circ}<\alpha<180^{\circ}$ nên $\cos\alpha<0$, vậy $\cos\alpha=\dfrac{-2\sqrt{2}}{3}$.\\ 
		 Khi đó $A=\cos\alpha+\sin\alpha=\dfrac{-2\sqrt{2}}{3}+\dfrac{1}{3}=\dfrac{1-2\sqrt{2}}{3}$.
	 }
\end{ex}

% nguồn: MaHoaNH_TracNghiem_Lop10
\begin{ex}%[0H4H1-3]%[ma:a1795fb2]
	Trong các đẳng  thức sau đây, đẳng thức nào \textbf{sai}?
	 \choice
		{$\sin  ^22x+\cos ^22x=1$}
		{$\cos \left(180^{\circ}-2x\right)=-\cos 2x$} 
		{$\sin  \left(90^{\circ}-x\right)=\cos x$}
		{\True $\sin ^2x+\cos ^22x=1$} 
	 \loigiai{
		$\sin  ^2x+\cos ^22x=1$ là đẳng thức sai.
	} 
\end{ex}

% nguồn: MaHoaNH_TracNghiem_Lop10
\begin{ex}%[0H4H1-3]%[ma:b77d6b43]
	Cho  $\triangle ABC$, hệ thức nào sau đây là đúng?
	\choice 
		{$\sin  (A+B)=\cos C$}
		{$\cos \left(\dfrac{A}{3} \right)=\sin \left(\dfrac{B+C}{3}  \right)$}
		 {$\sin (A+B+C.=1$}
		 {\True $\cos \left(\dfrac{A}{2} \right)=\sin \left(\dfrac{B+C}{2} \right)$}
	\loigiai{ 
		Ta  có $\sin \left(\dfrac{B+C}{2} \right)=\sin \left(\dfrac{180^\circ-A}{2} \right)=\sin \left(90^{\circ}-\dfrac{A}{2} \right)=\cos \left(\dfrac{A}{2} \right)$.
	} 
\end{ex}

% nguồn: MaHoaNH_TracNghiem_Lop10
\begin{ex}%[0H4B1-3]%[ma:bf575a20]
	Cho $2$ đường thẳng $d\colon x + y -2022=0$ và $d'\colon 2x - y + 2023=0$. Vị trí  tương đối của hai đường thẳng là
	\choice 
		{Vuông  góc}
		{\True  Cắt nhau}
		 {Song song}
		{Trùng nhau} 
	\loigiai{ Vì $\dfrac{1}{2}\ne \dfrac{1}{-1}$  nên hai đường thẳng $d$ và $d'$ cắt nhau
	 }
\end{ex}

% nguồn: MaHoaNH_TracNghiem_Lop10
\begin{ex}%[0H4Y1-3]%[ma:0e7a4cf8]
	Trong mặt phẳng  tọa độ $Oxy$, cho hai đường thẳng $d_1$ và $d_2$ có phương trình lần lượt là $d_1\colon a_1x+b_1y+c_1=0$ và $d_2\colon a_2x+b_2y+c_2=0$. Hệ phương trình $\heva{&a_1 x+b_1 y+c_1=0 \\ & a_2 x+b_2 y+c_2=0}$ có vô số nghiệm. Khi đó $d_1$ và $d_2$
	\choice 
		{song  song}
		 {cắt nhau nhưng không vuông góc nhau}
		{\True trùng nhau} 
		{vuông góc} 
	\loigiai{ 
		Hệ  phương trình vô số nghiệm khi $d_1$ và $d_2$ trùng nhau.
	 }
\end{ex}

% nguồn: MaHoaNH_TracNghiem_Lop10
\begin{ex}%[0H4Y1-3]%[ma:7f91bdd2]
	Trong mặt phẳng tọa độ $Oxy$, cho hai đường thẳng $d_1\colon 3x-4y+1=0$ và $d_2\colon x-y+2023=0$. Khẳng định nào  sau đây đúng?
	 \choice
		{$d_1$ song song với $d_2$} 
		 {$d_1$ trùng với $d_2$}
		{$d_1$  vuông góc với $d_2$}
		 {\True $d_1$ cắt $d_2$}
	 \loigiai{
		Do $\dfrac{3}{1}\ne  \dfrac{-4}{-1}$ nên hai đường thẳng $d_1$ và $d_2$ cắt nhau.
	 }
\end{ex}

% nguồn: MaHoaNH_TracNghiem_Lop10
\begin{ex}%[0H4B1-3]%[ma:be58cbaa]
	Phương trình nào sau đây là phương  trình đường thẳng không song song với đường thẳng $d \colon 3x-y-2=0$?
	 \choice
		{$-3x+y=0$} 
		 {$3x-y=0$}
		{$3x-y+6=0$} 
		 {$3x+y=0$}
	 \loigiai{
		Ta có $d$ không song song với $d' \colon  3x+y=0$ vì $\dfrac{3}{3} \ne \dfrac{-1}{1}$.
	 }
\end{ex}

% nguồn: MaHoaNH_TracNghiem_Lop10
\begin{ex}%[0H4K1-3]%[ma:b9b5a76c]
	Cho  đường thẳng $d_1 \colon 2x+3y+15=0$ và $d_2 \colon x-2y-3=0$. Khẳng định nào sau đây là đúng?
	 \choice
		{\True $d_1$ và $d_2$ cắt nhau  và không vuông góc với nhau}
		{$d_1$ và  $d_2$ song song với nhau}
		{$d_1$  và $d_2$ trùng nhau}
		{$d_1$ và $d_2$  vuông góc với nhau}
	\loigiai{ 
		Ta có $\dfrac{2}{1} \ne \dfrac{3}{-2}$ và $2  \cdot 1 +3 \cdot (-2) \ne 0$ nên $d_1$ và $d_2$ cắt nhau và không vuông góc với nhau,
	} 
\end{ex}

% nguồn: MaHoaNH_TracNghiem_Lop10
\begin{ex}%[0H4Y1-3]%[ma:e60c4e50]
	Trong mặt phẳng tọa độ, vị trí tương đối của hai  đường thẳng sau: $d_1: x-2 y+1=0$ và $d_2:-3 x+6 y-10=0$ là
	 \choice
		 {Song song}
		 {Cắt nhau nhưng không vuông góc nhau}
		 {Trùng nhau}
		{Vuông góc với  nhau}
	 \loigiai{
		Ta  có $ \dfrac{1}{-3}=\dfrac{-2}{6}\ne \dfrac{1}{-10} $.\\
		Do đó hai đường thẳng $ d_1 $ và $ d_2 $ song  song.
	} 
\end{ex}

% nguồn: MaHoaNH_TracNghiem_Lop10
\begin{ex}%[0H4B1-3]%[ma:3a2266e9]
	Trong mặt phẳng với hệ tọa độ $Oxy$, cho  đường thẳng $(\Delta)\colon 2x-y+1=0$. Đường thẳng nào sau đây song song với $(\Delta)$?
	\choice 
		{$2x+y-1=0$} 
		 {$x+y-1=0$}
		 {$x+2y-1=0$}
		{\True  $4x-2y+3=0$}
	 \loigiai{
		Đường thẳng $(\Delta)$  có một véc-tơ pháp tuyến là $\overrightarrow{n}=(2;-1)$.\\
		Suy  ra đường thẳng song song với $(\Delta)$ có dạng $2mx-my+c=0$, với $m\neq 0$, $c\neq m$.
	} 
\end{ex}

% nguồn: MaHoaNH_TracNghiem_Lop10
\begin{ex}%[0H4B1-3]%[ma:6fbef9a5]
	Tìm các giá trị thực của tham số $m$ để đường thẳng $y=\left(m^2-3\right)x+3m+1$ song song với đường  thẳng $y=x-5$.
	\choice 
		{$\hoac{ & m=2  \\ & m=-2}$}
		 {$m=-2$}
		 {\True $m=2$}
		{$\hoac{ & m=\sqrt{2} \\ &  m=-\sqrt{2}}$}
	 \loigiai{
	Đường  thẳng $y=\left(m^2-3\right)x+3m+1$ song song với đường thẳng $y=x-5$ khi
	$$\heva{ & m^2-3=1 \\ & 3m+1\neq-5} \Leftrightarrow \heva{  & \hoac{ & m=2 \\ & m=-2} \\ & m\neq-2} \Leftrightarrow m=2.$$
	} 
\end{ex}

% nguồn: MaHoaNH_TracNghiem_Lop10
\begin{ex}%[0H4K1-3]%[ma:e44b3383]
	Cho hai đường thẳng:  $2 x-y-1=0$ và $x+2 y+2=0$. Khi nói về vị trí tương đối của chúng, khẳng định nào đúng?
	 \choice
		{Cắt nhau nhưng không  vuông góc}
		{\True Vuông  góc}
		{Song song} 
		 {Trùng nhau}
	 \loigiai{
		Hai đường thẳng đã  cho có hai vectơ pháp tuyến lần lượt là $\vec{n}_1=(2;-1)$, $\vec{n}_2=(1;2)$.\\
		Vì $\vec{n}_1\cdot \vec{n}_2=2\cdot 1+(-1)\cdot 2=0$ nên  hai đường thẳng đã cho vuông góc với nhau.
	} 
\end{ex}

% nguồn: MaHoaNH_TracNghiem_Lop10
\begin{ex}%[0H4B1-3]%[ma:e3d63d4e]
	Cho đường thẳng $d_1\colon 2x+3y+15=0$ và $d_2\colon x-2y-3=0$. Khẳng định nào sau đây đúng? 
	 \choice
		 {$d_1$ và $d_2$ trùng nhau}
		{$d_1$ và  $d_2$ vuông góc với nhau}
		 {$d_1$ và $d_2$ song song với nhau}
		{\True $d_1$ và $d_2$ cắt nhau và không  vuông góc với nhau}
	 \loigiai{
		 Ta có $d_1\colon 2x+3y+15=0$ và $d_2\colon x-2y-3=0$ có véc-tơ pháp tuyến lần lượt là
		$$\vec{n}_{d_1}=(2;3) \text{ và  } \vec{n}_{d_2}=(1;-2).$$
		Do $\heva{&\dfrac{2}{1} \neq \dfrac{3}{-2}\\&\vec{n}_{d_1} \cdot \vec{n}_{d_2}=-4  \neq 0}$
		nên $d_1$ và $d_2$ cắt nhau và không vuông góc với  nhau.
	 }
\end{ex}


%%%%%%%%%%%%%%%% ĐÚNG/SAI %%%%%%%%%%%%%%%%

% nguồn: MaHoaNH_DungSai_Lop10
\begin{ex}%[0H4N1-3]%[ma:af6f0bae]
	Với mỗi góc $\alpha\ \left(0^\circ\leq  \alpha\leq 180^\circ\right)$. Các công thức dưới đây đúng hay sai?
	 \choiceTF
		{\True $\sin\alpha  =\tan\alpha\cdot\cos\alpha,\ \forall \alpha\neq 90^\circ$}
		{\True $\tan^2\alpha =1-\dfrac{1}{\cos^2\alpha},\ \forall  \alpha\neq 90^\circ$}
		 {$\cot^2\alpha =1+\dfrac{1}{\sin^2\alpha},\ \forall \alpha\in\left(0^\circ;180^\circ\right)$}
		{$\cos\alpha =\tan\alpha\cdot\sin\alpha  ,\ \forall \alpha\neq 90^\circ$}
	\loigiai{ 
		\begin{itemchoice}
			\itemch Ta có $\tan\alpha\cdot \cos\alpha =\dfrac{\sin\alpha}{\cos\alpha}\cdot\cos\alpha  =\sin\alpha,\ \forall \alpha\neq 90^\circ$.
			\itemch Ta có $\tan^2\alpha  =\dfrac{\sin^2\alpha}{\cos^2\alpha}=\dfrac{1-\cos^2\alpha}{\cos^2\alpha}=1-\dfrac{1}{\cos^2\alpha},\ \forall \alpha\neq 90^\circ$.
			\itemch Ta có $\cot^2\alpha =\dfrac{\cos^2\alpha}{\sin^2\alpha}=\dfrac{1-\sin^2\alpha}{\sin^2\alpha}=1-\dfrac{1}{\sin^2\alpha},\ \forall \alpha\in \left(0^\circ;180^\circ\right)$. 
			\itemch Ta có $\tan\alpha\cdot  \sin\alpha =\dfrac{\sin\alpha}{\cos\alpha}\cdot\sin\alpha =\dfrac{\sin^2\alpha}{\cos\alpha}=\dfrac{1-\cos^2\alpha}{\cos\alpha}=\dfrac{1}{\cos\alpha}-\cos\alpha,\ \forall \alpha\neq 90^\circ$.
		\end{itemchoice}
	} 
\end{ex}

% nguồn: MaHoaNH_DungSai_Lop10
\begin{ex}%[0H4H1-3]%[ma:9c19bea6]
	Mỗi công thức dưới đây đúng hay sai? 
	\choiceTF 
		{$\sin^2\alpha -\cos^2\alpha =1$} 
		{\True $\cos^4\alpha -\sin^4\alpha =2\cos^2\alpha -1$} 
		 {\True $\cos(180^\circ-\alpha)+\cos\alpha=\sin(180^\circ-\alpha)-\sin\alpha$}
		{\True $\sin^4\alpha +\cos^4\alpha =1-2\sin^2\alpha\cos^2\alpha$} 
	 \loigiai{
		\begin{itemchoice}
			\itemch Ta có $\sin^2\alpha -\cos^2\alpha  =\sin^2\alpha -\left(1-\sin^2\alpha \right)=1-2\sin^2\alpha \neq 1$.
			\itemch Ta có\\ $\begin{aligned}
			\cos^4\alpha -\sin^4\alpha &=\left(\cos^2\alpha-\sin^2\alpha\right)\left(\cos^2\alpha+\sin^2\alpha\right)\\ 
			 &=\cos^2\alpha-\sin^2\alpha\\
			&=\cos^2\alpha-\left(1-\cos^2\alpha\right)\\ 
			 &=2\cos^2\alpha -1.
		\end{aligned}$
		\itemch  Ta có \\
		$\text{VT} =\cos(180^\circ-\alpha)+\cos\alpha=-\cos\alpha+\cos\alpha=0$.\\ 
		$\text{VP}  =\sin(180^\circ-\alpha)-\sin\alpha=\sin\alpha-\sin\alpha=0$.
		\itemch Ta có $\tan\alpha\cdot \sin\alpha  =\dfrac{\sin\alpha}{\cos\alpha}\cdot\sin\alpha =\dfrac{\sin^2\alpha}{\cos\alpha}=\dfrac{1-\cos^2\alpha}{\cos\alpha}=\dfrac{1}{\cos\alpha}-\cos\alpha,\ \forall \alpha\neq 90^\circ$.
	\end{itemchoice}
} 
\end{ex}


%%%%%%%%%%%%%%%% TỰ LUẬN %%%%%%%%%%%%%%%%

% nguồn: MaHoaNH_TuLuan_Lop10
\begin{bt}%[0H4H1-3]%[ma:5f2b410c]
	Rút gọn biểu thức $A=\tan^2x\cdot\sin^2x-\tan^2x-\sin\left(90^\circ-x\right)\cdot\cos\left(180^\circ-x\right)$ ta được? 
	 \loigiai{
		Ta có 
		\allowdisplaybreaks 
		\begin{eqnarray*}
			A&=&\tan^2x\cdot\left(\sin^2x-1\right)-\cos x\cdot(-\cos  x)\\
			&=&\dfrac{\sin^2x}{\cos^2x}\cdot\left(-\cos^2x\right)+\cos^2x\\ 
			 &=&-\sin^2x+\cos^2x\\
			&=&\cos 2x. 
	\end{eqnarray*}}
\end{bt}

% nguồn: MaHoaNH_TuLuan_Lop10
\begin{bt}%[0H4H1-3]%[ma:c07f1c7f]
	Cho $\sin x+\cos x=\dfrac{2}{3}$, tính giá trị của biểu thức $P=\sin  ^3x+\cos ^3x+\sin x\cos x$.
	 \loigiai{
		Ta  có
		\allowdisplaybreaks 
		\begin{eqnarray*}
			 \sin x+\cos x=\dfrac{2}{3} &\Rightarrow & \left(\sin x+\cos x\right)^2=\dfrac{4}{9}\\
			&\Leftrightarrow &  \sin ^2x+\cos ^2x+2\sin x\cos x=\dfrac{4}{9}\\
			 &\Leftrightarrow & 2\sin x\cos x=-\dfrac{5}{9}\\
			&\Rightarrow &  \sin x\cos x=-\dfrac{5}{18}.
		\end{eqnarray*}
		 \[P=\left(\sin x+\cos x\right)^3-3\sin x\cos x\left(\sin x+\cos x\right)+\sin x\cos x=\dfrac{23}{27}-\dfrac{5}{18}=\dfrac{31}{54}.\]
	} 
\end{bt}

% nguồn: MaHoaNH_TuLuan_Lop10
\begin{bt}%[0H4H1-3]%[ma:13cfb23d]
	Cho $\tan \alpha=2$, tính giá trị của biểu thức $A=\dfrac{1+\cos  \alpha}{\sin \alpha} \left[1-\dfrac{\left(1-\cos \alpha \right)^2}{\sin ^2\alpha} \right]$.
	\loigiai{ 
		Ta có 
		\begin{align*}
			A&=\dfrac{1+\cos \alpha}{\sin  \alpha} \left[1-\dfrac{\left(1-\cos a\right)^2}{1-\cos ^2\alpha} \right]\\
			&=\dfrac{1+\cos \alpha}{\sin \alpha} \left[1-\dfrac{\left(1-\cos a\right)^2}{\left(1-\cos  \alpha \right)\left(1+\cos \alpha \right)}\right]\\
			&=\dfrac{1+\cos \alpha}{\sin  \alpha} \left[1-\dfrac{1-\cos a}{1+\cos \alpha} \right]\\
			&=\dfrac{1+\cos \alpha}{\sin \alpha} \left[\dfrac{1+\cos \alpha-\left(1-\cos a\right)}{1+\cos \alpha}  \right]\\
			&=\dfrac{1+\cos \alpha}{\sin \alpha}\cdot  \dfrac{2\cos \alpha}{1+\cos \alpha}=2\cot \alpha=\dfrac{2}{\tan \alpha}.
		\end{align*}
		$\Rightarrow  A=1$.
	 }
\end{bt}

% nguồn: MaHoaNH_TuLuan_Lop10
\begin{bt}%[0H4V1-3]%[ma:536afa47]
	Cho $\cot\alpha=-3$. Tính giá  trị biểu thức $P=\dfrac{\sin^3\alpha+\cos^3\alpha}{\sin\alpha-\cos\alpha}$.
	\loigiai{ 
		Vì $\cot\alpha=-3$ nên $\sin\alpha\ne 0$. Chia cả tử và mẫu  của $P$ cho $\sin^3\alpha$, ta có\\
	 \[P=\dfrac{1+\cot^3\alpha}{\dfrac{1}{\sin^2\alpha}(1-\cot\alpha)}=\dfrac{1+\cot^3\alpha}{\left(1+\cot^2\alpha\right)(1-\cot\alpha)}=\dfrac{1+(-3)^3}{\left[1+(-3)^2\right]\left[1-(-3)\right]}=-\dfrac{13}{20}=-0{,}65.\]}
\end{bt}

% nguồn: MaHoaNH_TuLuan_Lop10
\begin{bt}%[0H4V1-3]%[ma:2fee9ff3]
	Cho $\tan  \alpha=\dfrac{2}{5}$. Tính giá trị của biểu thức $A=\dfrac{4\sin \alpha + 3\cos \alpha}{5\sin \alpha - 2\cos \alpha}$.
	 \loigiai{
		Vì $\tan \alpha=\dfrac{\sin \alpha}{\cos  \alpha}=\dfrac{2}{5}$ nên $\cos \alpha \ne 0$.\\
		Chia cả tử và mẫu của $A$ cho $\cos \alpha$, ta được: 
		 \[
		A=\dfrac{4 \dfrac{\sin \alpha}{\cos \alpha}+3}{5\dfrac{\sin \alpha}{\cos \alpha}-2}=\dfrac{4\tan \alpha+3}{5\tan \alpha-2}=\dfrac{4\cdot \dfrac{2}{5}+3}{5\cdot \dfrac{2}{5}-2}= \dfrac{\dfrac{8}{5}+3}{2-2}=\dfrac{\dfrac{8  + 15}{5}}{0}= \dfrac{23}{0}=-\dfrac{23}{5}.
		 \]
	 }
\end{bt}

% nguồn: MaHoaNH_TuLuan_Lop10
\begin{bt}%[0H4V1-3]%[ma:9644951a]
	Tính giá trị biểu thức  $P=\sqrt{\sin^{4}x+6\cos^{2}x+3\cos^{4}x}+\sqrt{\cos^{4}x+6\sin^{2}x+3\sin^{4}x}$
	 \loigiai{
		Ta có $\sin^{4}x = (1-\cos^{2}x)^{2} = 1-2\cos^{2}x+\cos^{4}x$.\\ 
		$P_1  = \sin^{4}x+6\cos^{2}x+3\cos^{4}x = (1-2\cos^{2}x+\cos^{4}x)+6\cos^{2}x+3\cos^{4}x = 4\cos^{4}x+4\cos^{2}x+1 \\P_1= (2\cos^{2}x+1)^{2}$.\\
		$P_2 = \cos^{4}x+6\sin^{2}x+3\sin^{4}x = (1-2\sin^{2}x+\sin^{4}x)+6\sin^{2}x+3\sin^{4}x = 4\sin^{4}x+4\sin^{2}x+1 \\P_2=  (2\sin^{2}x+1)^{2}$.\\
		$P=\sqrt{(2\cos^{2}x+1)^{2}}+\sqrt{(2\sin^{2}x+1)^{2}}$.\\ 
		Vì $2\cos^{2}x+1>0$ và $2\sin^{2}x+1>0$  nên\\
		 $P=(2\cos^{2}x+1)+(2\sin^{2}x+1)=2(\cos^{2}x+\sin^{2}x)+2=2(1)+2=4$.
	} 
\end{bt}

% nguồn: MaHoaNH_TuLuan_Lop10
\begin{bt}%[0H4V1-3]%[ma:b1219f7c]
	Rút gọn biểu thức  $P=\dfrac{2\sin \left(180^{\circ}-x\right)\sin \left(90^{\circ}-x\right)}{\tan x+\cot x}+\sin ^4x+\cos ^4x$.
	 \loigiai{
		Ta  có
		 \allowdisplaybreaks
		\begin{eqnarray*}
			P&=&\dfrac{2\sin  x\cos x}{\dfrac{\sin x}{\cos x}+\dfrac{\cos x}{\sin x}}+\left(\sin ^2x+\cos ^2x\right)^2-2\sin ^2x\cos ^2x\\
			&=&\dfrac{2\sin x\cos x}{\dfrac{1}{\sin x\cos x}}+1-2\sin ^2x\cos ^2x\\ 
			 &=&2\sin ^2x\cos ^2x+1-2\sin ^2x\cos ^2x\\
			&=&1. 
		\end{eqnarray*}
	} 
\end{bt}

% nguồn: MaHoaNH_TuLuan_Lop10
\begin{bt}%[0H4B1-3]%[ma:ce3415b4]
	Xét vị trí tương đối của hai đường thẳng 
	 $d_1\colon\heva{&x=-2+t\\ &y=3-2t}$ và $d_2\colon\heva{&x=3-2t'\\ &y=-1+4t'.}$
	 \loigiai{
		Ta có $d_1$ có véc-tơ chỉ phương $\overrightarrow{u_1}(1;-2)$  và $d_2$ có véc-tơ chỉ phương $\overrightarrow{u_2}(-2; 4)$.\\ Suy ra $\overrightarrow{u_2}=-2\overrightarrow{u_1}$ nên $\overrightarrow{u_1},\overrightarrow{u_2}$ cùng phương. Do đó, chúng song song hoặc trùng nhau.\\
		Mặt khác $M(-2; 3)\in  d_1$ nhưng $M\notin d_2$ nên $d_1\parallel d_2$.
	 }
\end{bt}
